\section{Experiments}
\label{sec:experiments}

\subsection{Benchmark graphs and preprocessing}
The selected benchmark graphs are drawn from the Harwell--Boeing collection.
The raw matrix files and normalized graph files are stored separately. Each
entry in the versioned manifest records the source page, raw-file path, graph
path, matrix type, number of vertices, number of normalized undirected edges,
maximum degree, number of isolated vertices, and SHA-256 checksums for both
raw and normalized data. A further checksum is computed from the canonical
sequence consisting of $n$ and the sorted edge list.

For a square matrix $A$, vertices correspond to rows and columns. An undirected
edge $\{i,j\}$ is included for $i\ne j$ if at least one of $A_{ij}$ or
$A_{ji}$ is nonzero. Diagonal entries are discarded, duplicate and
opposite-direction entries are merged, and isolated vertices are retained.
Consequently, the number of nonzero matrix entries is not generally the number
of edges in the normalized simple graph. For example, the source reports 50
stored matrix entries for JGL009~\cite{NISTJGL009}, whereas the normalized graph used here has 32 unique
undirected edges. We therefore do not copy the paper's matrix-entry counts into
the normalized manifest.

\begin{table}[t]
\centering
\caption{Three Harwell--Boeing graphs used in the 300-second main benchmark.}
\label{tab:dataset}
\begin{tabular}{lrrrrr}
\hline
Instance & $n$ & $m$ & $\Delta(G)$ & Isolates & Density \\
\hline
JGL009 & 9 & 32 & 8 & 0 & 0.8889 \\
BCSSTK01 & 48 & 176 & 11 & 0 & 0.1560 \\
BP 400 & 822 & 4015 & 294 & 0 & 0.0119 \\
\hline
\end{tabular}
\end{table}

JGL009 and BCSSTK01 also occur in Table A.2 of Ren et al.~\cite{Ren2020CBMP}
and their source records are available from NIST~\cite{NISTJGL009,NISTBCSSTK01}.
BP 400 is an additional Harwell--Boeing matrix from the SMTAPE set~\cite{NISTBP400}
and is used as a larger sparse instance. BCSSTM01 is documented separately as
an excluded case: after removing its diagonal entries, it has 48 isolated
vertices and no edges~\cite{NISTBCSSTM01}. The two distributed forms of JGL009
(plain and gzip-compressed) are the same matrix and are counted once. The
repository also records the 28 names in
Table A.2 as candidates; a candidate is included in a benchmark only
after its raw file, normalized graph, and checksums have been verified.

\subsection{Pilot protocol and observed outcomes}
The day-8 pilot covers three verified graphs and the five methods available in the original run: pairwise SAT, order SAT, CP-SAT, CPLEX MIP, and
Gurobi MIP. Each planned run has a 60-second end-to-end wall-clock budget,
one thread, and seed zero. The original run log contains 15 records. Table~\ref{tab:day8-pilot}
uses the 60-second CPLEX Studio 22.2 reruns on JGL009 and BCSSTK01
and the 60-second Gurobi 13.0.3 rerun on all three graphs. In the table, OPT means
optimality proven, F/TO means a validated feasible incumbent with timeout, and
Guard means rejected by the preflight size guard before solver startup.

On JGL009 the identity bandwidth and the degree lower bound both equal 4. Both SAT runs stopped before CNF construction. On BCSSTK01,
OR-Tools CP-SAT proved bandwidth 12 optimal in 46.686 s (record 8 of \path{results/pilot_day8.jsonl}). In the earlier day-6 pilot, pairwise SAT and CP-SAT each reached 12 with lower bound 6. Both SAT encodings return
validated feasible solutions but do not prove optimality within 60 seconds.
On BP 400, CP-SAT returns a feasible solution but does not prove optimality
within the time limit; the pairwise and order SAT encodings are rejected by
the preflight guard because their estimated clause counts are approximately
3.27 billion and 2.99 billion, respectively. The guard adds $mn^2$ possible edge clauses to each assignment count before choosing $K$. This upper estimate includes every ordered position pair; solver time for a guarded run is N/A.

The new 60-second CPLEX Studio reruns proved bandwidth 4 optimal on
JGL009 and returned a validated feasible layout of bandwidth 24 on
BCSSTK01, with lower bound 6 and timeout. Thus the latter entry has
bounds $[6,24]$ and carries no optimality proof. The Gurobi academic-license rerun proved bandwidth 4 optimal
on JGL009. On BCSSTK01 it returned bounds $[6,17]$ at timeout. BP~400 exceeded
the runner's size guard before either MIP solver started. Both MIP solvers
also returned validated layouts on the smaller graphs under the separate
300-second protocol (Table~\ref{tab:six-method-rerun}).

The five-method pilot excludes IBM CP Optimizer. The separate six-method
60-second pilot applies a 1 GiB Windows job memory cap. Gurobi exhausted the memory allowance on BCSSTK01 before returning a layout. The independent
order/CP pilot is identified in the Order Encoding section.

\subsection{Main benchmark protocol}
The six-method rerun used an AMD Ryzen 5 5600H with 15.4 GiB of
OS-visible RAM, Windows 11 Home (build 26300), and Python 3.12.3.
Each instance--method pair had a 300-second target including model construction,
one requested solver thread, and seed zero where supported. Cleanup and
native solver return can add elapsed time beyond the target. The six-method
60-second pilot limited total Windows job committed memory to 1024 MiB
for the runner and its children. The 300-second run used model-size guards
without that memory cap. Some runs overlapped other local work; elapsed
times here describe these runs and are unsuitable for a speed ranking.

The benchmark runner stores one JSON object per run in JSONL format and
supports resuming completed runs. It preserves failures and distinguishes
solver status (optimal, feasible, or no solution) from termination reason
(completed, timeout, out-of-memory, license error, or other error). Every
reported incumbent is independently validated. The five-method, 60-second pilot and the six-method, 300-second run are
separate experiments. The pilot outcomes are reported above; the six-method
batch is summarized in the section on MIP and six-method results.
